\section{Типичные ошибки и их устранение}

В процессе работы с дроном часто возникают ситуации, когда <<код правильный, но ничего не работает>>. В этой главе собраны самые частые проблемы.

\subsection{Проблемы с камерой и маркерами}

\subsubsection{Дрон не видит маркер}
\begin{itemize}
    \item \textbf{Причина 1: Плохое освещение.} Камере темно, выдержка увеличивается, появляется <<смаз>> (Motion Blur) при движении.
    \item \textbf{Решение:} Включите дополнительный свет. Маркер должен быть освещен равномерно.
    
    \item \textbf{Причина 2: Блики.} Маркер заламинирован глянцевой пленкой или лежит под лампой.
    \item \textbf{Решение:} Переместите маркер или используйте матовый скотч.
    
    \item \textbf{Причина 3: Неверный размер.} В коде указан размер 0.2 метра, а напечатан маркер 18 см.
    \item \textbf{Решение:} Измерьте черную рамку линейкой и исправьте константу в коде.
    
    \item \textbf{Причина 4: Quiet Zone.} Маркер вырезан по черному контуру без белого поля.
    \item \textbf{Решение:} Наклейте его на белый лист бумаги.
\end{itemize}

\subsubsection{Дрон летит не туда}
\begin{itemize}
    \item \textbf{Причина 1: Перепутаны оси.} Вы думаете, что X — это вправо, а это вперед.
    \item \textbf{Решение:} Проверьте ориентацию (Body Frame). Ось X всегда смотрит <<из камеры>> (вперед).
    
    \item \textbf{Причина 2: Ошибка калибровки.} Камера <<рыбий глаз>> сильно искажает картинку по краям.
    \item \textbf{Решение:} Проведите калибровку камеры заново.
    
    \item \textbf{Причина 3: Yaw (Рысканье).} Дрон повернулся на 90 градусов, но код продолжает отправлять команды в старой системе координат.
    \item \textbf{Решение:} Используйте глобальную систему координат (World Frame) или учитывайте текущий угол Yaw.
\end{itemize}

\subsection{Проблемы с подключением и Wi-Fi}

\subsubsection{Connection Refused / Timeout}
\begin{itemize}
    \item \textbf{Причина 1:} Компьютер подключился к другой Wi-Fi сети (с интернетом).
    \item \textbf{Решение:} Проверьте, что вы подключены к точке доступа Pioneer\_XXXX.
    
    \item \textbf{Причина 2:} Брандмауэр Windows блокирует порт 8000/8888.
    \item \textbf{Решение:} Разрешите Python доступ в частные сети.
\end{itemize}

\subsection{Ошибки в коде Python}

\subsubsection{IndentationError}
Python очень чувствителен к отступам. Не смешивайте табы и пробелы. Используйте 4 пробела для отступа.

\subsubsection{OpenCV Error: Assertion failed}
Обычно означает, что функция получила пустое изображение (None).
\begin{itemize}
    \item Проверьте, что \texttt{camera.get\_frame()} не вернул None.
    \item Проверьте, что файл изображения существует (если читаете с диска).
\end{itemize}

\subsection{Аварийные ситуации}

\subsubsection{Дрон не реагирует на Land}
Если программная посадка не срабатывает:
\begin{enumerate}
    \item Не паникуйте.
    \item Нажмите кнопку Disarm в Ground Station (если есть).
    \item Аккуратно поймайте дрон снизу за аккумулятор (берегите пальцы от пропеллеров!) и переверните его вверх ногами. Моторы отключатся автоматически.
\end{enumerate}

\begin{warnbox}{Внимание!}
Никогда не пытайтесь схватить дрон сверху или сбоку — пропеллеры вращаются со скоростью 10 000 об/мин и могут нанести серьезные травмы.
\end{warnbox}
